The environment milestone is substantially complete, and a working baseline has undergone initial validation. The project is transitioning from the proposed weeks 5-6 baseline milestone toward weeks 7-8, which focus on strategic biases. Final rule variations and style definitions were deferred rather than completed during the originally proposed weeks 1-2.

The next task is to document a fixed baseline and design the four strategy-specific biases. A proposed approach is to add small configurable adjustments to candidate scores while keeping the cutoff heuristic and rollout policy shared. Behavioral metrics must establish that rear-priority, front-priority, aggressive, and defensive agents actually differ as intended. Equal search settings will be maintained, with measured runtime reported because candidate counts can differ.

After those pilots, rule variations will be selected around explicit strategic incentives and prewritten hypotheses. Each variation should change a clearly identified mechanism, and experiments should measure how often it affects play. The proposed scope is two core variations initially, with additional variations only if time permits. The exact choices remain undecided.

Later work will run balanced tournaments under standard and modified rules, collecting win rates, game lengths, hits, pip-count progression, exposed blots, and rule-activation frequency. Results will be compared against standard-rule baselines and accompanied by uncertainty estimates. Pilot seeds will be separated from final evaluation seeds. Final deliverables remain the documented simulator and agents, gameplay datasets, analysis, report, and presentation.

Key remaining challenges are heuristic-weight justification, rollout cost, preservation of experiment records, and separating strategic style from playing strength. Optional comparison with GNU Backgammon remains future work. Conclusions will be limited to the implemented agents and tested variations rather than universal claims about human Backgammon strategy.
